\textbf{Время:} 01:00

\textbf{Тезис:} импульсы измеряются, а фемтометровые точки испускания требуют косвенного исследования.

\textbf{На экране:} неизвестная область слева → стрелки импульсов → измеряемое и неизвестное справа → нижняя фраза о корреляциях.

\textbf{Текст:} На предыдущем слайде мы выделили область испускания. Теперь слева она показана вопросительным знаком, а стрелки обозначают два конечных импульса. Над областью стоят наши вопросы: размер и время испускания.

Справа сформулировано различие: импульсы известны — их восстанавливают по траекториям частиц. А фемтометровые координаты испускания недоступны разрешению детектора. Один фемтометр — десять в минус пятнадцатой степени метра.

Поэтому нужен косвенный способ. Нижняя фраза предлагает искать информацию о размере в корреляциях: не только в каждом импульсе отдельно, но и в том, как часто два импульса встречаются вместе.

\textbf{Переход:} Увидим, как похожую задачу о размере решали при наблюдении звезды.

\textbf{Справка:} Общая вершина столкновения измеряется; неизвестны именно индивидуальные фемтометровые точки и времена испускания. Стрелки — схема известных импульсов, а не восстановленные траектории внутри источника. Устройство TPC и TOF объясняется позднее, в экспериментальном блоке.

\textbf{Источники:} Lisa et al. (2005), введение и §3.
